\chapter{Listas e Dicionários}

\section{Listas}

\subsection{Criação}

\begin{lstlisting}
let nums  = [1, 2, 3, 4, 5]
let nomes = ["RS", "SP", "RJ"]
let misto = [1, "dois", true]
let vazia = []
\end{lstlisting}

\subsection{Acesso por índice}

Índices começam em 0. Índices negativos contam a partir do fim:

\begin{lstlisting}
let xs = [10, 20, 30, 40]
display xs[0]     // 10
display xs[3]     // 40
display xs[-1]    // 40  (último)
display xs[-2]    // 30
\end{lstlisting}

Acesso fora dos limites produz erro de runtime:

\begin{lstlisting}
xs[10]
// error: index 10 out of bounds (len=4)
\end{lstlisting}

\subsection{Comprimento}

\begin{lstlisting}
len([1, 2, 3])    // 3
\end{lstlisting}

\subsection{Mutação}

\lstinline{push} adiciona ao final da lista (muta no lugar, retorna \lstinline{nil}):

\begin{lstlisting}
let lista = [1, 2, 3]
push(lista, 4)
display lista    // [1, 2, 3, 4]
\end{lstlisting}

\begin{aviso}
  \lstinline{push} é uma função mutante — altera \lstinline{lista} diretamente e
  retorna \lstinline{nil}. Não faça \lstinline{let nova = push(lista, 4)}.
\end{aviso}

\lstinline{pop} remove e retorna o último elemento:

\begin{lstlisting}
let ultimo = pop(lista)
\end{lstlisting}

\subsection{Slicing}

\begin{lstlisting}
let xs = [0, 1, 2, 3, 4, 5]
display slice(xs, 1, 4)    // [1, 2, 3]  (índices 1..3)
\end{lstlisting}

\subsection{Iteração}

\begin{lstlisting}
for x in [10, 20, 30] {
  display x * 2
}
\end{lstlisting}

\subsection{Funções funcionais}

\begin{lstlisting}
let xs = [1, 2, 3, 4, 5]

let dobrados  = map(xs, fn(x) { x * 2 })
let pares     = filter(xs, fn(x) { x % 2 == 0 })
let soma      = reduce(xs, 0, fn(acc, x) { acc + x })
\end{lstlisting}

\section{Ranges}
\label{sec:ranges}

Ranges são expressões que produzem listas de inteiros. A sintaxe segue Rust:

\begin{lstlisting}[caption={Ranges como listas}]
let a = 1..5      // [1, 2, 3, 4]   (exclusivo)
let b = 1..=5     // [1, 2, 3, 4, 5] (inclusivo)
\end{lstlisting}

Por serem listas normais, ranges funcionam com qualquer função que aceite lista:

\begin{lstlisting}
map(1..=5, fn(x) { x^2 })            // [1, 4, 9, 16, 25]
filter(1..10, fn(x) { x % 2 == 0 }) // [2, 4, 6, 8]
len(1..100)                          // 99
\end{lstlisting}

\subsection{\texttt{chain}}

\lstinline{chain} concatena múltiplas listas (ou ranges) em uma só:

\begin{lstlisting}[caption={chain com ranges e listas}]
chain(1..3, 9..=11)           // [1, 2, 9, 10, 11]
chain(1..=3, [10, 20], 5..7)  // [1, 2, 3, 10, 20, 5, 6]
\end{lstlisting}

Uso típico em laços sobre sequências não-contíguas:

\begin{lstlisting}
for lag in chain(1..=3, 6..=8) {
  display f"lag {lag}"
}
\end{lstlisting}

\subsection{Composição com closures}

Ranges, \lstinline{chain} e closures compõem naturalmente em uma única expressão.
A sintaxe \lstinline{|x| expr} define uma closure sem a cerimônia de \lstinline{fn}:

\begin{lstlisting}[caption={Expressão composta — única linha}]
let r = map(chain(1..=10, 5..9, 56..70), |n| n * 3)
\end{lstlisting}

O pipe também funciona no cabeçalho do \lstinline{for}, compondo transformações
antes de entrar no corpo do laço:

\begin{lstlisting}[caption={Pipeline no cabeçalho do for}]
for i in chain(1..=10, 5..9, 56..70) |> filter(|n| (n % 2) == 0) {
    display i
}
\end{lstlisting}

\subsection{Comparação com outras linguagens}

A tabela abaixo mostra como as principais linguagens expressariam o mesmo pipeline:

\begin{center}
\begin{tabular}{ll}
\toprule
\textbf{Linguagem} & \textbf{Forma} \\
\midrule
Python  & \texttt{filter(lambda n: n\%2==0, chain(range(1,11), range(5,9), range(56,70)))} \\
Rust    & \texttt{(1..=10).chain(5..9).chain(56..70).filter(|n| n\%2==0)} \\
Haskell & \texttt{filter even \$ [1..10] ++ [5..8] ++ [56..69]} \\
Julia   & \texttt{Iterators.filter(n->n\%2==0, Iterators.flatten([1:10,5:8,56:69]))} \\
R       & \texttt{Filter(function(n) n\%\%2==0, c(1:10, 5:8, 56:69))} \\
\midrule
Hayashi & \texttt{chain(1..=10,5..9,56..70) |> filter(|n| (n\%2)==0)} \\
\bottomrule
\end{tabular}
\end{center}

\hayashi{} é a única linguagem em que o pipeline é estritamente da esquerda para a
direita, sem imports, sem namespaces e com \lstinline{chain} variádico que trata
todos os argumentos simetricamente.

\subsection{Nota sobre desempenho}

\lstinline{chain} materializa os ranges em uma lista antes de iterar.
\lstinline{for i in 1..N}, por outro lado, itera diretamente sem alocar nada.

\lstinline{chain} é para sequências não-contíguas pequenas — lags, anos, grupos de
variáveis. Esse é o seu único caso de uso legítimo. Qualquer sequência contígua é
simplesmente \lstinline{1..N} ou \lstinline{1..=N}.

\section{Dicionários}

\subsection{Criação}

\begin{lstlisting}
let config = {
  "host": "localhost",
  "porta": 5432,
  "ssl": true
}

let vazio = {}
\end{lstlisting}

Chaves podem ser strings ou inteiros. Valores podem ser de qualquer tipo.

\subsection{Acesso}

\begin{lstlisting}
config["host"]     // "localhost"
config["porta"]    // 5432
\end{lstlisting}

Chave ausente produz erro:

\begin{lstlisting}
config["usuario"]
// error: key 'usuario' not found in dict
\end{lstlisting}

Use \lstinline{try}/\lstinline{catch} para acesso seguro:

\begin{lstlisting}
let user = try { config["usuario"] } catch _ { "anonimo" }
\end{lstlisting}

\subsection{Inserção e atualização}

\begin{lstlisting}
let d = {"a": 1}
let d2 = dict_set(d, "b", 2)    // d inalterado, d2 = {"a":1,"b":2}
\end{lstlisting}

\lstinline{dict_set} é funcional — retorna um novo dicionário sem mutar o original.

\subsection{Verificando chave}

\begin{lstlisting}
"host" in config    // true
"senha" in config   // false
\end{lstlisting}

\subsection{Dicionários como retorno de funções}

Funções de análise frequentemente retornam dicionários com múltiplos resultados:

\begin{lstlisting}[caption={summarize retorna dicionário}]
load "dados.csv" as df
let stats = summarize(df, renda)

display stats["mean"]    // média
display stats["sd"]      // desvio padrão
display stats["N"]       // número de observações
\end{lstlisting}

\subsection{Iteração sobre dicionários}

\begin{lstlisting}
let d = {"a": 1, "b": 2, "c": 3}

for key in keys(d) {
  display f"{key} = {d[key]}"
}
\end{lstlisting}
